\minitopic{比较不是寻找相同词}两个传统都有可译为“知”的词，不表示它们提出同一问题。词的意义取决于与哪些词对立、能进入哪些推理、在哪些实践中使用。现代分析知识常与信念、真和确证联系；古典文本中的“知”还可能与学、行、智、名、道等构成网络。 反过来，没有对应词也不表示没有相关思想。传统可能通过学习制度、辩论规范和修养实践处理错误与可靠性，而不把它抽象成“认识论”。比较单位可从词扩大到论证、案例和实践，但必须说明转换步骤。 同源词内部也变化。先秦、宋明与现代汉语的“知”不可默认恒定；英文 knowledge 在不同时期也有理论负担。比较研究同时需要历史语义与哲学重构。

\minitopic{三层写作规范}第一层写\textbf{文本事实}：某篇、某章、某角色说了什么，版本和翻译如何。第二层写\textbf{解释主张}：这段话在全篇与历史语境中承担什么功能。第三层写\textbf{哲学重构}：若把它转成当代问题，可形成怎样的论证。三层应使用不同语气。 “《庄子》写道……”需要文本定位；“这段批评固定视角”是解释；“它可发展成一种语境认识论”是重构。若省略层次，读者会把研究者的创造性方案误当历史作者明说的理论。 重构并非不合法。哲学传统一直在新问题中再解释旧资源。规范是公开目标、保留文本阻力，并允许重构因反例而修正，而不是借古人权威为现代主张背书。

\minitopic{《论语》中的知、不知与学习}《论语》把承认不知、好学、听言观行和识人等放在同一实践中\parencite{analects2003}。这提供认识谦逊、证言判断和长期能力培养的资源，却不是一套JTB定义。核心单位常是学习者怎样成为能判断和行动的人。 “知之为知之，不知为不知”可解释为二阶校准\parencite[2.17]{analects2003}：主体不仅形成一阶判断，还应准确判断自己是否拥有知识。它反对虚假自信，但不说明一阶知识全部条件。放到AI界面中，可成为系统表达不确定性的规范性类比，而不是历史预见。 孔子重视多闻、多见与择取，也强调言行一致。证言信任不是无条件服从：观察行为和长期记录构成对言说者的评价。师生关系同时包含权威与学习自主，适合与第九章的关系性自主对话。

\minitopic{孟子的反身认识}孟子讨论心、性、恻隐等道德能力，反身工夫旨在扩充和实现道德端绪\parencite{mencius2008}。把它直接等同于笛卡尔内省，会忽略对象与目标：前者不是从不可怀疑心理表象奠定外部世界知识。 可比问题是第一人称对道德动机是否有特殊入口，情感经验能否成为道德理由，以及修养如何改善自知。孟子也承认环境与习惯影响能力，这使“内在”不等于脱离社会因果。 现代重构可提出：某些规范理由通过受教育的情感显现。反对者会问情感如何区分文化偏见，支持者需加入反思、他人批评和实践结果。这样的对话产生新论证，而非简单贴“内在主义”标签。

\minitopic{荀子的解蔽与学习}荀子强调人的认识能力、对象之理、语言分类和心受“蔽”限制，并把学习、师法与辩说视为矫正条件\parencite{xunzi2014}。蔽不是一次感官幻觉，而是执著一端遮蔽整体。它与确认偏误、立场局限有结构相似。 解蔽不要求没有视角，而要求通过比较、秩序和公共规范扩展判断。这里可与社会认识论的批评制度对话。差异在于荀子的“道”、礼与政治秩序具有实质规范内容，不是中立的信息多样性原则。 若把“兼陈万物而中县衡”重构为多源证据和权衡方法，应说明这是现代抽象\parencite[chap.~21]{xunzi2014}。文本对师法和正名的重视还带来权威问题：规范秩序能校正偏见，也可能压制异议。比较应保留这项双重性。

\minitopic{庄子的视角与相对主义}《庄子》频繁转换人物、物种与梦醒视角，揭示判断标准和语言区分的局限\parencite{zhuangzi2009}。把它归为“凡观点都同样真”的简单相对主义不充分，因为文本也批评僵化、展示技艺和适应情境的优劣。 与当代语境主义的相似在于标准与视角重要；差异在于语境主义通常仍给知识归属确定真值条件，庄子更广泛地松动执着于固定是非和自我中心。二者的问题范围不同。 庄周梦蝶可与梦怀疑比较，但笛卡尔用梦寻找确定基础，庄子让身份与转化边界不再自明。相同故事装置服务相反方向。比较价值正来自这一差异，而非谁更早提出怀疑论。

\minitopic{墨家、名辩与理由规范}墨家文本重视“故”、类、名实与辩，提供公共论证和分类的材料。将其直接套成亚里士多德三段论会损失文本结构；更适合比较理由如何支持主张、类比何时成立、名称怎样对应实践区分\parencite{graham1989}。 “白马非马”等名辩命题涉及名称范围、类与复合限定。现代内涵—外延工具可帮助重构，却不是证明古代作者已经拥有Frege式语义学。重构应展示多个解释及其文本代价。 公共可争辩性是重要连接点。理由若只能由私人体验把握，难以进入辩；名、类和标准使判断可共享。现代形式逻辑提供更严格符号工具，但同样依赖选取恰当分类和前提。

\minitopic{佛教量论与二谛的比较边界}佛教传统内部差异巨大。因明与量论讨论现量、比量、可靠认知和推理标志；中观二谛讨论世俗语言与胜义分析。把所有内容合成“佛教认识论”会抹去学派、语言和历史层次\parencite{dunne2004,siderits2013}。 现量不能简单等于当代感觉材料，比量也不只是西方演绎。它们嵌在错误认知、概念构造和解脱目标中。与知觉、推理理论比较前，需要确定具体作者与论证。 “世俗谛=现象、胜义谛=物自体”的康德式对应尤其危险。两组概念在形而上学和方法中功能不同\parencite{siderits2013}。可比较两层话语怎样约束真理主张，但不能用术语配对代替解释。

\minitopic{工夫论对德性认识论的挑战}德性认识论常把能力归属于个体；工夫论强调身体习惯、情感、礼仪、师友和具体实践共同塑造判断。它推动我们把认知德性看作发展过程和关系成果，而非主体已有模块。 工夫也提供训练细节：反复实践、在事上检验、通过榜样和纠错形成敏感性。这比只列“开放、认真、谦逊”德性清单更动态。现代教育可借此问，制度怎样让德性可练习。 但工夫目标常是道德—政治人格，而非最大化真信念。若现代重构只保留提高准确率，会抽空传统价值。更好的对话让真理导向与人格、关系目标彼此提出问题。

\minitopic{关系性认识自主}西方现代图景常把自主想成独立于他人判断。证言与社会认识论已经显示，成熟自主是管理依赖。儒家师友与共同学习可提供关系性模型：他人帮助主体形成判断能力，而不只是替主体决定。 关系也可能压迫。年龄、身份和礼制权威若免于质疑，会损害纠错。关系性自主必须包含权威可问责、学习者逐步取得发言能力和转换关系的可能。 比较结论因而不是“儒家重关系更好”，而是提出评价问题：依赖是否培养能力，能否回应反证，退出与申诉是否可行，信用如何分配。这些标准可用于古代实践与现代教育。

\minitopic{去中心化与共同标准}扩大课程传统不是按地域各放一章就完成。若所有理论问题仍由单一传统定义，其他材料只作装饰，结构仍中心化。问题本身也应被改变：从“知识由哪些条件构成”扩展到“什么实践培养可知者”“认识目标为何”。 去中心化又不能取消准确性标准。文本需有版本与语境，论证需有效，比较需可批评。以尊重多元为由免于证据要求，会把非西方传统当不可讨论的身份象征，而非平等哲学伙伴。 课程可使用双向压力测试：用当代分析工具要求古典重构精确，用古典问题结构揭露当代理论的个体主义或命题中心倾向。双方都能改变，才是真正对话。

\minitopic{综合案例：如何比较“知行合一”}第一步做文本工作：说明具体版本、上下文和“知”“行”的用法。第二步列历史解释：命题是否针对把道德认可与实践割裂\parencite{tiwald2014}。第三步才作重构：完成的道德知包含相应行动倾向。 随后选择现代对话对象。与能力之知比较，问会做是否可还原为命题；与德性认识论比较，问稳定品格如何使真信念成为成就；与动机理论比较，问真诚道德判断是否必然推动行动。每次比较只承诺局部结构相似。 最后列出不可化约差异：良知的形上与伦理背景、工夫目标、现代知识理论的真值中心。一个合格比较结论应同时写“它帮助我们看见什么”和“它不能被翻译成什么”。

\begin{tcolorbox}[breakable,colback=white,colframe=blue!30!black,title={比较研究工作表},fonttitle=\bfseries]
记录原文、版本、说话角色与语境；列关键词的语义范围和替代译法；区分文本事实、历史解释、结构比较与规范重构；说明比较单位与目的；写出相似点、功能差异和不可化约处；让双方理论各自接受一个新问题或反例。
\end{tcolorbox}
